\documentclass[10pt]{article}

% --- Page Setup ---
\usepackage[margin=0.5cm]{geometry}
\usepackage{multicol}

% \allowdisplaybreaks

% --- Math and Code ---
\usepackage{amsmath, amssymb}
\usepackage{listings}

% --- Spacing Adjustments ---
\usepackage{enumitem}
\usepackage{titlesec}

% 1. Remove space between list items
\setlist{nosep, leftmargin=*} 

% 2. Shrink space around section headers
\titlespacing*{\section}{0pt}{2pt}{1pt}
\titlespacing*{\subsection}{0pt}{1ex plus .2ex minus .2ex}{0.5ex plus .1ex}

% 3. Add a thin line between columns for readability
\setlength{\columnseprule}{0.4pt}
\setlength{\columnsep}{1.5em}

% --- Custom Code Formatting ---
\lstset{
    basicstyle=\ttfamily\scriptsize,
    breaklines=true,
    aboveskip=2pt,
    belowskip=2pt
}

% \usepackage{xcolor}
% \newcommand{\sectionbox}[1]{%
%      \colorbox{blue!75!black}{\color{white}\small\bfseries~#1~}%
%    }
%    \titleformat{\section}{}{}{0em}{\sectionbox}
%    \titleformat*{\subsection}{\normalsize\bfseries}

\setlength{\parskip}{0pt}
\setlength{\baselineskip}{0pt plus 0.2pt} 

\begin{document}
\small 

\raggedright
% \raggedbottom

\begin{multicols*}{2}
\section*{Stastical ML}
\textbf{Supervised learning dataset}
$\mathcal{D} = \{(x_1, y_1), (x_2, y_2), \dots, (x_n, y_n)\}$\\
\textbf{Hypothesis function} $h : \mathcal{X} \mapsto \hat{y}$\\
\textbf{Binary classification} $\mathcal{Y} = \{+1,-1\}$.\\
\textbf{Multiclass classification} $\mathcal{Y} = \{1,2\ldots,K\},K>2$.\\
\vspace{1em}
\textbf{Probabilistic models}\\
\textbf{Sigmoid}\\
$\sigma(z) = \frac{1}{1 + \exp(-z)}$\\
\begin{itemize}
    \item $z\gg 0 \implies \sigma(z) \approx 1$
    \item $z \ll 0 \implies \sigma(z) \approx 0$
    \item $z = 0 \implies \sigma(z) = 0.5$
\end{itemize}
\textbf{Logistic regression (binary classification)}
$p(y = +1 \mid x, w) = \sigma(h_w(x)), \qquad p(y = -1 \mid x, w) = 1 - \sigma(h_w(x))$\\
Unified expression:
$p(y \mid x, w) = \frac{1}{1 + \exp(-y \cdot h_w(x))}$\\
Decision boundary: 
$\hat{y} = \begin{cases} +1 & \text{if } h_w(x) \geq 0 \\ -1 & \text{if } h_w(x) < 0 \end{cases}$\\
\textbf{MLE objective for log reg}
$\hat{w} = \arg\min\limits_{w \in \mathbb{R}^d} \frac{1}{n} \sum_{i=1}^{n} \log\!\left(1 + \exp(-y_i \cdot h_w(x_i))\right)$\\
From taking negative log-likelihood of the prob model and minimising it over the training data.
There is no closed-form solution, must ust gradient descent or other optimisers.\\
\medskip
\textbf{Softmax (multiclass)}\\
$q_{v,k} = \frac{\exp(v_k)}{\sum_{j=1}^{K} \exp(v_j)}, \qquad k = 1, \dots, K$\\
Takes a vector of $K$ scores and converts them into a valid prob distribution. The ranking of scores is preserved.\\
\textbf{Multiclass prob model}\\
$p(y \mid x, w) = \frac{\exp(h_{w,y}(x))}{\sum_{k=1}^{K} \exp(h_{w,k}(x))}$\\
This is softmax applied to $h_w(x)$ evaluated a $y$, the true class.\\
\textbf{Multiclass prediction rule}\\
$\hat{y} = \arg\max\limits_{y \in \mathcal{Y}} \; h_{w,y}(x)$\\
Softmax not needed at prediction time, the class with highest logit has the highest probability.
Softmax is only needed during training to compute probabilities for the loss function.\\
\textbf{Multiclass MLE objective (cross-entropy loss)}\\
$\min_{w} \; \frac{1}{n} \sum_{i=1}^{n} \underbrace{\left( -h_{w,y_i}(x_i) + \log \sum_{k=1}^{K} \exp(h_{w,k}(x_i)) \right)}_{\ell(h_w(x_i),\, y_i)}$\\
\begin{itemize}
    \item $-h_{w,y_i}(x_i)$pushes the true class score up
    \item $\log \sum_{k=1}^{K} \exp(h_{w,k}(x_i))$ grows whenever any score grows, acting as a penalty that keeps all other scores in check.
\end{itemize}
The loss is minimised when the true class score is large relative to all others — not just large in absolute terms.\\
\textbf{Pipeline}
\begin{itemize}
    \item Feed $x$ into $h_w(x) \in \mathbb{R}^K$ to get one score per $K$ classes.
    \item Apply softmax to get a valid probability distribution.
    \item Pick out the probability of the true class $y_i$ and calc the cross entropy loss (take -ve log).
    \item Update $w$ by gradiend descent (using gradient of the loss $\nabla_w l$).
\end{itemize}
\vspace{1em}
Loss functions are typically negative log-likelihoods of the probabilistic models.\\
\textbf{Logistic loss (binary)}\\
$\ell(y', y) = \log(1 + \exp(-y \cdot y'))$\\
$y$ is the true label, $y' = h_w(x)$ is the raw score.\\
When $yy'$ positive (prediction is correct and label agrees), the loss is small. When -ve, the loss grows. The loss never reaches 0, so the weights keep updating.\\
Logistic loss $\iff$ logistic/sigmoid model.\\
\textbf{Hinge loss (binary)}\\
$\ell(y', y) = \max\{1 - y \cdot y', \; 0\}$\\
Zero loss once $yy'\geq 1$ and stops updating weights.\\
\textbf{Exponential loss (binary)}\\
$\ell(y', y) = \exp(-y \cdot y')$\\
Penalises wrong predictions far more aggressively than logistic loss.
This makes it very sensitive to outliers and mislabelled data, since a single bad example can dominate the total loss.\\
\textbf{Unifying intuition for binary losses}
$\text{margin} = y \cdot y'
\begin{cases} > 0 & \text{correct prediction} \\ < 0 & \text{incorrect prediction} \end{cases}$\\
\textbf{Absolute loss (regression)}\\
$\ell(y', y) = |y - y'|$\\
Absolute loss $\iff$ Laplace model.\\
\textbf{Squared loss (regression)}\\
$\ell(y', y) = \frac{1}{2}(y - y')^2$\\
Standard loss for regression problems.
Penalises large errors more heavily due to the square.\\
Squared loss $\iff$ Gaussian model.\\
$\epsilon$\textbf{-sensitive loss (regression)}\\ 
$\ell(y', y) = \max\{|y - y'| - \varepsilon, \; 0\}$
Ignores errors smaller than $\epsilon$.
Only errors that fall outside the tube are penalised, and penalised linearly rather than quadratically, making it more robust to outliers than squared loss.\\
\textbf{Cross-entropy loss (multiclass)}
$\ell(h_w(x), y) = -h_{w,y}(x) + \log \sum_{k=1}^{K} \exp(h_{w,k}(x))$\\
Standard loss for multiclass classification with softmax.
This is equivalent to $-\log p(y | x, w)$, the negative log probability of the true class under the softmax model.\\
Cross-entropy model $\iff$ softmax model.\\
\textbf{KL divergence}
$D_{\text{KL}}(p \| q) = \sum_{k=1}^{K} p_k \log \frac{p_k}{q_k}$\\
Measures discrepancies between two prob distributions $p$ and $q$. Always $\geq 0$, and equals zero only when $p=q$.\\
\textbf{Cross-entropy loss $\equiv$ KL divergence minimisation}\\
$L(w) = D_{\text{KL}}\!\left(P[Y \mid X] \;\|\; p(\cdot \mid X, w)\right) + \underbrace{H(P[Y \mid X])}_{\text{constant w.r.t. } w}$\\
Minimising the expected cross-entropy loss is exactly equivalent to minimising the KL divergence between the true distribution $P[Y|X]$ and the model's predicted distribution $p(\cdot | x, w)$.\\
\vspace{1em}
\textbf{Expected loss}
$L(w) = \mathbb{E}_{(X,Y)}\left[\ell(h_w(X), Y)\right]$\\
This is the true objective: average loss over all possible inputs weighted by how likely they occur.\\
We can never compute it directly because we don't have the true data distribution, only a finite sample.\\
\textbf{Empirical loss}
$L_n(w) = \frac{1}{n} \sum_{i=1}^{n} \ell(h_w(x_i), y_i)$\\
The tractable approximation of $L(w)$ using only the training data. As $n\to\infty, L_n(w)\to L(w)$ by LLN.\\
But for finite $n$, minimising $L_n(w)$ does not guarantee minimising $L(w)$, this gap is the source of overfitting.\\
\textbf{ERM}
$\hat{w} = \arg\min\limits_{w} \; L_n(w)$
The standard training procedure — find the $w$ that minimises the average loss on training data.
Needs to be paired with regularisation or model selection to control overfitting.\\
\textbf{Overfitting}
$\left| \mathbb{E}_{(X,Y)}\left[\ell(\hat{h}(X), Y)\right] - \frac{1}{n}\sum_{i=1}^{n} \ell(\hat{h}(x_i), y_i) \right|$\\
Overfitting is when this gap is large — the model performs well on training data but poorly on new data.
Happens when the model is too complex relative to the amount of training data, allowing it to memorise the training set rather than learn the underlying pattern.
Most dangerous when data is scarce.\\
\textbf{Regularised objective}
$\min_{w} \; L_{n,\lambda}(w) = \frac{1}{n} \sum_{i=1}^{n} \ell(h_w(x_i), y_i) + \frac{\lambda}{2} \|w\|_2^2$\\
Adds a penalty on the size of $w$ directly to the training objective.
Forces a tradeoff — the model must fit the data well and keep weights small.\\
\begin{itemize}
    \item $\lambda$ too small: overfitting, regularisation no effect.
    \item $\lambda$ too large: underfitting, model too simple.
\end{itemize}
$\lambda$ is a hyperparameter chosen via cross-validation.\\
\textbf{L2 regularisation (Ridge)}
$\|w\|_2^2 = \sum_{j=1}^{p} w_j^2$\\
Penalises large weights by squaring them. Shrinks all weights uniformly toward zero but never exactly to zero — all features remain in the model.
Good default when you expect all features to contribute something.\\
\textbf{L1 regularisation (Lasso)}
$\|w\|_1 = \sum_{j=1}^{p} |w_j|$\\
Penalises weights by absolute value. Tends to push weights to exactly zero, effectively removing irrelevant features — produces sparse solutions.
Use when you expect only a few features to matter.\\
\textbf{Elastic Net}
$\alpha \|w\|_1 + \frac{1-\alpha}{2} \|w\|_2^2, \qquad \alpha \in (0,1)$\\
Combines L1 and L2. $\alpha$ closer to 1 gives more sparsity, $\alpha$ closer to 0 gives more uniform shrinkage.\\
\vspace{1em}
\textbf{Model class}
$\mathcal{H} = \{h_w : \mathcal{X} \to \mathcal{Y} \mid w \in \Omega \subset \mathbb{R}^p\}$\\
Rather than searching over all possible functions, you restrict the search to a specific parametric family.
The choice of $\mathcal{H}$ is made before training, determines what the model can represent.
A model $\mathcal{G}$ is more complex than $\mathcal{H}$ if $\mathcal{H}\subset \mathcal{G}$.\\
\textbf{Bias-variance decomp}
$L(w) \approx \text{Bias}^2 + \text{Variance}$\\
High bias = model too simple, cannot capture the true pattern regardless of how much data you have.\\
High variance = model too complex, changes drastically with different training sets.\\
\textbf{Holdout method}
$\mathcal{D} = \mathcal{D}_{\text{train}} \cup \mathcal{D}_{\text{val}} \cup \mathcal{D}_{\text{test}}$\\
\begin{enumerate}
    \item Train on $\mathcal{D}_{\text{train}}$.
    \item Evaluate on $\mathcal{D}_{\text{val}}$ to choose hyperparameters.
    \item Final evaluation on $\mathcal{D}_{\text{test}}$ once only.
\end{enumerate}
Validation loss is an unbiased estimate of expected loss because the model never sees validation data during training.
High variance: estimate depends heavily on which examples ended up in $\mathcal{D}_{\text{val}}$.\\
\textbf{$k$-fold cross validation}\\
$\text{CV score} = \frac{1}{k} \sum_{i=1}^{k} \text{validation loss on fold } i$\\
Divide data into $k$ equal subsets. For each fold $i$, train on the remaining $k - 1$ folds, evaluate on fold $i$. Average the $k$ validation losses.\\
Every example is used for both training and validation, just never at the same time.
More data-efficient and lower variance than holdout.
More expensive, requires $k$ separate training times.\\
Choose the hyperparameter value with the lowest CV score.\\
\textbf{Leave-one-out cross-validation (LOOCV)}\\
Special case of $k$-fold where $k=n$, so each fold has exactly one example.\\
Almost all data used for training each time — very low bias in the estimate. But requires $n$ separate training runs, computationally expensive.\\
Use only when data is very scarce and you cannot afford to leave more out.\\





\section*{Optimisation}
\textbf{Empirical Loss Minimisation (ERM)}\\
$\min_w \left\{ \frac{1}{n} \sum_{i=1}^{n} \ell(h_w(x_i), y_i) \right\}$\\
Minimises avg loss over $n$ training examples.\\
\textbf{Regularised ERM}\\
$\min_w \left\{ \frac{1}{n} \sum_{i=1}^{n} \ell(h_w(x_i), y_i) + \frac{\lambda}{2}\|w\|^2 \right\}$\\
Adds penalty to prevent overfitting. $\lambda$ as regularisation strength.
\textbf{MLE}\\
$\max_{\theta \in \Theta} \left\{ \sum_{i=1}^{n} \log p_\theta(X_i) \right\}$\\
Finds parameters that make observed data most probable. Equivalent to minimising cross-entropy loss (ERM).\\
\textbf{MAP}\\
$\max_{\theta \in \Theta} \left\{ \sum_{i=1}^{n} \log p(X_i|\theta) + \log p(\theta) \right\}$\\
MLE + prior belief $\log p(\theta)$ about parameters.\\
Gaussian prior $p(\theta) = \mathcal{N}(0,\sigma^2)$ gives $\log p(\theta) \propto -\|\theta\|^2$, making MAP identical to regularised ERM.\\
\textbf{Connections (optimisation $\iff$ probabilistic)}\\
ERM $\iff$ MLE; regularised ERM $\iff$ MAP; L2 penalty $\frac{\lambda}{2}\|w\|^2 \iff$ Gaussian prior.\\[1\baselineskip]

\subsubsection*{Mathemtical properties of functions}
\textbf{Lipschitz Smoothness}\\
$\|\nabla\phi(w) - \nabla\phi(w')\| \leq L\|w - w'\|, \quad \forall w, w' \in \mathbb{R}^d$\\
Gradient cannot change abruptly between any two points.\\
Small $L$: gradient changes slowly; Large $L$: gradient changes faster but still controlled.\\
Sufficient condition: all eigenvalues of Hessian $\nabla^2\phi(w) \preceq LI_d$ everywhere.\\
\textbf{Lipschitz key inequality}\\
$|\phi(w') - \phi(w) - \nabla\phi(w)^\top(w'-w)| \leq \frac{L}{2}\|w'-w\|^2$\\
Error of tangent approximation is at most $\frac{L}{2}\|w' - w\|^2$.\\
Gives quadratic upper bound on function - used to prove each GD step decreases objective.\\
\textbf{Strong convexity}\\
$\phi(w') \geq \phi(w) + \nabla\phi(w)^\top(w'-w) + \frac{\mu}{2}\|w'-w\|^2, \quad \forall w, w' \in \mathbb{R}^d$
Gives quadratic lower bound on function.\\
Guarantees unique global minimum, no flat regions or saddle points.\\
Sufficient condition: all eigenvalues of Hessian $\nabla^2\phi(w) \succeq \mu I_d$ everywhere.\\
\textbf{PL-Condition (from strong convexity)}\\
$\|\nabla\phi(w)\|^2 \geq 2\mu\left(\phi(w) - \min_w \phi(w)\right)$\\
Gradient norm is always large enough relative to suboptimality gap.\\
\begin{itemize}
    \item Far from minimum: gradient is large, GD takes big steps.
    \item Close to minimum: gradient is small but never vanishes before reaching minimum.
\end{itemize}
Guarantees convergence of GD.\\
\textbf{Smoothness + Strong Convexity Together}\\
$\mu I_d \preceq \nabla^2\phi(w) \preceq LI_d$\\
Function is sandwiched between two parabolas with curvatures $\mu$ and $L$.\\
\textbf{Condition number} $\kappa = \frac{L}{\mu}$\\
$\kappa \approx 1$: nearly spherical landscape, easy to optimise.\\
$\kappa \gg 1$: elongated narrow landscape, GD zigzags, many iterations needed.

\subsubsection*{GD}
\textbf{GD update rule} $w_{t+1} = w_t - \eta_t\nabla F(w_t)$\\
$\eta_t$: step size at step $t$.\\
Moving in negative gradient direction always locally decreases objective:\\
$\frac{d}{dt}F(w - t\nabla F(w))\bigg|_{t=0} = -\|\nabla F(w)\|^2 < 0$\\
\textbf{GD step size} $\eta = \frac{1}{L}$\\
Guaranteed decrease in objective function $F(w_t)$ per step.\\
Larger $L$: smaller steps, more cautious of overshooting.\\
Too large $\eta > \frac{2}{L}$: overshoots, objective can increase.\\
Too small: too slow.\\
\textbf{Guaranteed decrease per step} (from Lipschitz smoothness)\\
$F(w_{t+1}) \leq F(w_t) - \frac{1}{2L}\|\nabla F(w_t)\|^2$\\
Each step decreases objective by at least $\frac{1}{2L}\|\nabla F(w_t)\|^2$.\\
Larger $L$: larger changes in gradient, step size is smaller, smaller guaranteed decrease per step.\\
\textbf{Exponential shrinking of gap} (from PL-condition)\\
$F(w_T) - F(w_*) \leq \left(1 - \frac{\mu}{L}\right)^T(F(w_0) - F(w_*))$\\
LHS is the suboptimality gap at step $T$, how far we are from the optimal loss. RHS is upper bound on how large the gap can be.\\
$0 < (1-\frac{\mu}{L}) < 1$, so repeated multiplication of this factor shrinks the gap exponentially toward 0.\\
\begin{itemize}
    \item $\kappa \approx 1$: $(1-\frac{\mu}{L}) \approx 0$, gap shrinks fast per step.
    \item $\kappa \gg 1$: factor close to 1, many steps needed to close the gap.
\end{itemize}
To guarantee the gap $< \epsilon$, solve for $T$ s.t. RHS $<\epsilon$, giving $T \geq \frac{L}{\mu}\log\frac{F(w_0)-F(w_*)}{\epsilon}$.\\
\textbf{GD iteration complexity}
$O\left(\kappa\log\frac{1}{\epsilon}\right)$\\
$\propto\kappa$: large condition number, slow convergence.\\
$\log \frac{1}{\epsilon}$: 10x more accuracy requires fixed no. of extra steps.\\[1\baselineskip]

\textbf{Scaling problem with GD}\\
\textbf{Finite sum structure}
$\min\limits_{w \in \mathbb{R}^d} \left\{ F(w) = \frac{1}{n}\sum_{i=1}^{n} f_i(w) \right\}$\\
$f_i(w)$: loss on single training example $i$.\\
Computing $\nabla F(w) = \frac{1}{n} \sum_{i=1}^{n} \nabla f_i(w)$ requires accessing every $\nabla f_i$ per step (a full pass through dataset every iteration).\\
We need to consider a second complexity beyond iteration complexity: \textbf{total complexity}, total number of $\nabla f_i$ accesses to achieve $\epsilon$-error.\\

\begin{tabular}{|c|c|c|}
\hline
& Iteration Complexity & Total Complexity \\
\hline
GD & $O\left(\kappa\log\frac{1}{\epsilon}\right)$ & $O\left(n\kappa\log\frac{1}{\epsilon}\right)$ \\
\hline
\end{tabular}
Iteration complexity independent of $n$ - only counts steps.\\
Total complexity = iteration complexity $\times$ n, since each step accesses all $n$ gradients.\\
Scales directly with dataset size - fundamental problem with GD on large datasets.\\[\baselineskip]
\textbf{SGD}\\
\textbf{SGD update rule}
$w_{t+1} = w_t - \eta_t \nabla f_{i_t}(w_t), \quad i_t \sim \text{Uniform}\{1, 2, \dots, n\}$\\
Randomly picks one data point $i_t$ per step instead of computing the full gradient over all $n$.\\
Each step is $n$ times cheaper than GD.\\
\textbf{Unbiasedness of SGD}
$\mathbb{E}_{i_t}[\nabla f_{i_t}(w_t)] = \nabla F(w_t)$\\
Stochastic gradient is valid substitute for full gradient: on average equals the true gradient.\\
Any single step may be noisy, but SGD heads in the right direction on average.\\
\textbf{SGD step size}
$\eta = \min\left\{\frac{\mu\epsilon}{L\sigma^2}, \frac{1}{L}\right\}$\\
Same $\frac{1}{L}$ term as GD step size, but also consider $\frac{\mu\epsilon}{L\sigma^2}$ to control gradient noise. If the steps are noisier, we need smaller steps.\\
% \begin{tabular}{|c|c|c|c|c|}
% \hline
% & \multicolumn{2}{c|}{Iteration Complexity} & \multicolumn{2}{c|}{Total Complexity} \\
% \hline
% & GD & SGD & GD & SGD \\
% \hline
% Strongly convex 
% & $O\left(\kappa\log\frac{1}{\epsilon}\right)$ 
% & $\tilde{O}\left(\kappa\log\frac{1}{\epsilon} + \frac{\kappa\sigma^2}{\mu\epsilon}\right)$ 
% & $O\left(n\kappa\log\frac{1}{\epsilon}\right)$ 
% & $\tilde{O}\left(\kappa\log\frac{1}{\epsilon} + \frac{\kappa\sigma^2}{\mu\epsilon}\right)$ \\
% \hline
% \end{tabular} 

\begin{tabular}{|c|c|c|}
\hline
& Iteration Complexity & Total Complexity \\
\hline
SGD & $\tilde{O}\left(\kappa\log\frac{1}{\epsilon} + \frac{\kappa\sigma^2}{\mu\epsilon}\right)$ & $\tilde{O}\left(\kappa\log\frac{1}{\epsilon} + \frac{\kappa\sigma^2}{\mu\epsilon}\right)$ \\
\hline
\end{tabular}
Iteration complexity (SGD) has extra term compared to GD: $\frac{\kappa\sigma^2}{\mu\epsilon}$: additional steps from gradient noise $\sigma^2$ - how much individual $\nabla f_i$ differ from each other.\\
SGD total complexity is free from $n$: each step only accesses 1 gradient.\\
For large $n$, SGD is better than GD on total complexity.\\
$\tilde{O}$: same as $O$ but hides additional log factors.\\[\baselineskip]

\section*{Sampling}
\textbf{Monte-Carlo Integration}\\
Goal: estimate an expectation that is analytically intractable:\\
$\mathbb{E}_{X \sim p}[f(X)] = \int f(x)p(x)dx$\\
\begin{itemize}
    \item Values of $f(x)$ weighted by how probable each $x$ is under $p$.
    \item Intractable when $p$ is too complex or integral is over very high dimensional space.
\end{itemize}
\textbf{Monte-Carlo Approximation}
$\frac{1}{M}\sum_{i=1}^{M} f(x_i) \approx \mathbb{E}_{X \sim p}[f(X)], \quad x_1, \dots, x_M \sim p(x)$\\
Draw $M$ i.i.d. samples from $p$, average $f$ over them.\\
By LLN, approximation error shrinks to 0 as $M \to \infty$.\\
Error rate $O(\frac{1}{\sqrt{M}})$: dimension-free, unlike numerical integration.\\
\textbf{Why not standard numerical integration?}\\
\begin{itemize}
    \item Curse of dimensionality: computationally impossible for high dimensions.
    \item MC error $O(\frac{1}{\sqrt{M}})$ does not depend on dimension at all.
\end{itemize}
\textbf{The fundamental challenge}\\
MC requires drawing samples from $p(x)$, which is often hard because:\\
\begin{itemize}
    \item $p(x)$ may only be known up to a normalising constant $Z$: computing $Z = \int \tilde{p}(x)dx$ is itself intractable.
    \item $p(x)$ may be a complex high dimensional distribution like a posterior in Bayesian inference.
\end{itemize}
\vspace{1em}
\textbf{Recap - Bayesian Regression}\\
Posterior:
$p(w|\mathbf{y}) \propto p(\mathbf{y}|w, \mathbf{X})p(w) = \prod_{i=1}^{n} p(y_i|x_i, w)p(w)$\\
Predictive distribution:\\
$p(y|\mathbf{y}, x) = \int p(y|w, x)p(w|\mathbf{y})dw$\\
\textbf{MC approximation of predictive}:
$\frac{1}{M}\sum_{i=1}^{M} p(y|w_i, x) \approx p(y|\mathbf{y}, x), \quad w_1, \dots, w_M \sim p(w|\mathbf{y})$\\
Draw samples from posterior, average predictions over them.\\
Requires sampling from $p(w|\mathbf{y})$: intractable directly.\\[\baselineskip]

\textbf{LMC}\\
Update rule:
$w_{t+1} = w_t + \eta\nabla\log p(w_t) + \sqrt{2\eta}\xi_k, \quad \xi_k \sim \mathcal{N}(0, I_d)$\\
\begin{itemize}
    \item $\eta\nabla\log p(w_t)$: gradient term, moves toward regions of high probability.
    \item $\sqrt{2\eta}\xi_k$: noise term, random Gaussian pertubation added every step.
\end{itemize}
\textbf{LMC vs MC integration}\\
MC Integration: the goal — approximate intractable integral by averaging over samples. Assumes samples already available.\\
LMC: the mechanism — generates samples from complex $p(x)$ when direct sampling impossible.\\
Together: LMC generates samples $\to$ MC integration uses them to approximate integral.\\
\textbf{2 issues with LMC}\\
Burn-in period: Early particles start far from target distribution - not representative samples.
Discard initial phase of chain, only retain samples after convergence.\\
Sample correlation: Consecutive $w_t, w_{t+1}$ highly correlated - each step moves a small distance.
Solution is thinning - retain every $k$-th sample, discard intermediate ones to reduce correlation.\\
\textbf{Bayesian Inference w/ LMC}\\
LMC update for Posterior:\\
$w_{t+1} = w_t + \eta\nabla\log p(w_t|\mathbf{y}) + \sqrt{2\eta}\xi_k, \quad \xi_k \sim \mathcal{N}(0, I_d)$\\
where the gradient of the log posterior is:\\
$\nabla\log p(w_t|\mathbf{y}) = \sum_{i=1}^{N}\nabla\log p(y_i|x_i, w) + \nabla\log p(w)$\\
\begin{itemize}
    \item $\sum_{i=1}^{N}\nabla\log p(y_i|x_i, w)$ pushes $w$ toward parameters that explain data well.
    \item $\nabla\log p(w)$ prior term, for Gaussian prior $\mathcal{N}(0,\sigma^2)$ this equals $-\frac{w}{\sigma_0^2}$, pulling $w$ back toward zero.
\end{itemize}
Same tension as regularised ERM - balancing data fit and prior belief.\\
\textbf{Full pipeline}\\
$\underbrace{\text{LMC}}_{\text{sample from } p(w|\mathbf{y})} \rightarrow \underbrace{w_1, \dots, w_M}_{\text{after burn-in + thinning}} \rightarrow \underbrace{\frac{1}{M}\sum_{i=1}^{M}p(y|w_i, x)}_{\text{MC approximation of } p(y|\mathbf{y},x)}$\\
\textbf{Connection between optimisation and sampling}\\
The likelihood gradient $\sum_{i=1}^{N}\nabla\log p(y_i|x_i, w)$ has the same finite sum structure as ERM.\\[0.5\baselineskip]

\begin{tabular}{|c|c|c|}
\hline
& \textbf{Optimization} & \textbf{Sampling} \\[6pt]
\hline
Method & GD / SGD & LMC \\[6pt]
\hline
Update & $-\eta\nabla F(w)$ & $+\eta\nabla\log p(w) + \sqrt{2\eta}\xi_k$ \\[6pt]
\hline
Output & Single point $w^*$ & Distribution of samples \\[6pt]
\hline
Finite sum & $\dfrac{1}{n}\displaystyle\sum_{i=1}^{n}\nabla f_i$ & $\displaystyle\sum_{i=1}^{n}\nabla\log p(y_i|x_i,w)$ \\[12pt]
\hline
Scaling problem & GD $\rightarrow$ SGD & LMC $\rightarrow$ SG-LMC* \\[6pt]
\hline
\end{tabular}
*Stochastic-gradient LMC not in lecture.

% \begin{tabular}{|c|c|c|}
% \hline
% & \textbf{Optimization} & \textbf{Sampling} \\[6pt]
% \hline
% Method & GD / SGD & LMC \\[6pt]
% \hline
% Update & $-\eta\nabla F(w)$ & $+\eta\nabla\log p(w) + \sqrt{2\eta}\xi_k$ \\[6pt]
% \hline
% Output & Single point $w^*$ & Distribution of samples \\[6pt]
% \hline
% Finite sum & $\dfrac{1}{n}\displaystyle\sum_{i=1}^{n}\nabla f_i$ & $\displaystyle\sum_{i=1}^{n}\nabla\log p(y_i|x_i,w)$ \\[12pt]
% \hline
% \end{tabular}

\end{multicols*}
\end{document}